\documentclass{article}
\usepackage{amsmath,hyperref}
\begin{document}
\section*{Native radio receipts need a reusable resource-context boundary}
Use the separation in the
\href{https://github.com/eternal-november/eos-research-programmes/blob/4359a8a1484498b9c46d8e58a0112448a1fba06b/projects/contextual-composition/publications/interactads/Interactads.tex}{Interactads presentation}:
one must say what the actual resource context retains when a
qualified content receipt is compared or replayed. The present
scope is a fixed native radio accounting period, registered
participant identities, nonempty valid frames and known test
links. Payload outcomes are content observations; the physical
rate context is indexed independently by sender/receiver, direction
and unexpired native debit time. No proof of all relative
pseudodistributive-law hypotheses is asserted.

\paragraph{Projection, future checks, and scope.}
Group the exact native charged rows for each scoped endpoint and
direction by debit time. The within-one-second use curve at
future cut $t$ is the sum of lengths at stamp $s\ge t-1$.
Rows that are older can be dropped; $s=t-1$ still counts.
A content retry, receipt read or equality of current balances
is not a native resource update. At a fixed endpoint with every
possible byte demand available to probe and no other bottleneck,
two different future use curves differ in at least one future
admission test: at that cut choose the maximum packet fitting
the less used resource. Under that declared broader probe scope,
a current balance cannot replace the time-indexed profile.
The executable groups same-time histories with different packet
partitions and checks their future consequences. It does not
claim equivalence of old frame audit records when those are
an admissible observable.

\paragraph{A direct pasting test.}
At one cut, with each named link working, to deliver all intended
frames it is necessary and sufficient that the summed new charges
per actual sender and per actual receiver fit their respective
remaining native budgets. Necessity follows from native accepted
appends; sufficiency follows by the prefix charge bound under
any order if the entire demand fits. Separate tests for different
named packets need not suffice because they may name the same
physical endpoint/direction. The source program verifies
15,000 such batch fixtures, 42,013 native attempts including
earlier debit or erasure, and 3,276 cases where both isolated
checks pass but the shared cut cannot perform both. The projected
state is verified after every actual native attempt, not inferred
from an aggregate success count.

For non-successful attempts, native send-side admission can charge
before link or receive-side rejection. Therefore the content
branch of a failed request is insufficient to prove that the
effect was free. Within this finite tested grammar, two
equivalent time-indexed profiles have the same successor profile
and public delivery outcome given the same next typed request
and compatible link context. This gives a concrete operational
descent argument for ordered continuations. A live endpoint
replacement may preserve or change the resource identity; no
silent clearing of its prior debits is assumed. This transports
the native credit cases to explicit compatibility and profile
proof obligations, not a novel radio model or an independent
mission result.
\end{document}
